\section*{Chapter \ref{ch:iv:python} --- Python}

\begin{solution}{iq:iv:python:1}
\texttt{[1]}, \texttt{[1, 2]}, \texttt{[3]}. The default list is created once, when \texttt{def} runs, and every call without the argument
appends to the same list; the third call passes its own. Fix: \texttt{def add\_fill(fill, fills=None)} and \texttt{if fills is None:
fills = []}. A linter flags the pattern (ruff's B006), which is how to find it without being told.
\iqlookfor{evaluation time of defaults, the standard fix, and tooling that catches it.}
\end{solution}

\begin{solution}{iq:iv:python:2}
\texttt{[20, 20, 20]} then \texttt{[0, 10, 20]}. The lambdas capture the variable \texttt{k}, not its value, and are called after the loop
ends with \texttt{k = 2}. A default argument \texttt{k=k} is evaluated when each lambda is created, capturing the current value;
\texttt{functools.partial} does the same.
\iqlookfor{late binding of closures and the default-argument fix.}
\end{solution}

\begin{solution}{iq:iv:python:3}
\texttt{True True} then \texttt{True False}. The language guarantees only the equalities. CPython keeps one object for each integer from
$-5$ to 256, so the two 256s are the same object, while \texttt{int('1000')} creates a new object each time. Never use \texttt{is} to
compare numbers or strings; use it for \texttt{None} and sentinels.
\iqlookfor{equality against identity, and naming the CPython detail as a detail.}
\end{solution}

\begin{solution}{iq:iv:python:4}
\texttt{5} then \texttt{0}: the first \texttt{sum} consumes the generator ($0 + 1 + 4$) and the second finds it exhausted. Use a list if the
values are needed twice, or recreate the generator.
\iqlookfor{single-pass iterators.}
\end{solution}

\begin{solution}{iq:iv:python:5}
Deleting a key while iterating over the dictionary changes its size, and CPython raises \texttt{RuntimeError: dictionary changed size
during iteration} at the next step. Build a new dictionary with a comprehension (the last line), or iterate over a copy of the keys
(\texttt{for sym in list(positions)}).
\iqlookfor{the mutation-during-iteration rule and two safe patterns.}
\end{solution}

\begin{solution}{iq:iv:python:6}
Stream it: a generator that reads one line at a time, a generator that filters, and a dictionary accumulating notional per symbol
(\cref{lst:iv:python:pipeline}); memory is the dictionary plus one line, whatever the file's size. The chapter's test runs it on
5\,000 generated fills and checks it against a direct computation. For speed rather than memory, read in chunks with pandas or
Polars and aggregate per chunk.
\iqlookfor{constant-memory streaming with generators.}
\end{solution}

\begin{solution}{iq:iv:python:7}
\omcode{code/interviews/24-python/python/iv_py.py}{28}{36}{Restore the limit in a \texttt{finally}
clause.}\label{lst:iv:python:limit}
The \texttt{finally} runs whether the block ends normally or by an exception, and the exception still propagates. The chapter's test
checks both paths.
\iqlookfor{\texttt{contextlib.contextmanager} with \texttt{try}/\texttt{finally}.}
\end{solution}

\begin{solution}{iq:iv:python:8}
\texttt{[100. 0. 102. 103.]} and \texttt{True False}. \texttt{prices[1:3]} is a view, so writing to it changes \texttt{prices[1]};
\texttt{prices[[2, 3]]} is a copy, so writing to it changes nothing in \texttt{prices}. The function
\texttt{np.shares\_memory} confirms which is which.
\iqlookfor{basic slicing gives views, fancy indexing gives copies.}
\end{solution}

\begin{solution}{iq:iv:python:9}
\texttt{False 0.9999999999999999} and \texttt{True}: 0.1 is not exactly representable in binary, and the running sum accumulates the
rounding; \texttt{math.fsum} tracks the lost low-order bits and returns the correctly rounded sum. On a desk it matters when a
position or cash balance is checked for exact zero, when many small fills are summed, and when two systems reconcile totals computed
in different orders: keep money in integer units (cents, ticks) and compare floats with a tolerance.
\iqlookfor{binary representation, compensated summation, and integer units for money.}
\end{solution}

\begin{solution}{iq:iv:python:10}
The work is mostly waiting on sockets, so \texttt{asyncio} (one thread, 300 streams, cheap task switches) or a thread per group of
streams both fit; threads help because they release the lock while blocked. If the per-message statistic is heavy pure Python,
the lock serialises it: move it to numpy (which releases the lock), to a compiled extension, or to worker processes fed through
queues. The lock does not make compound operations atomic: a shared counter or dictionary updated from several threads still needs a
lock (or a queue to a single owner).
\iqlookfor{matching the concurrency model to I/O-bound and CPU-bound work, and the limits of the GIL.}
\end{solution}

\begin{solution}{iq:iv:python:11}
With $c$ the cumulative sum of the prices preceded by a zero, the mean of the window ending at $i$ is $(c_{i+1} - c_{i+1-k})/k$
(\cref{lst:iv:python:rolling}). Check it against the obvious loop on random inputs of random lengths and windows, as the chapter's
test does for 50 seeds; for long series of large numbers, beware the cancellation in differences of large cumulative sums, and use
\texttt{np.convolve} or pandas' \texttt{rolling} for stability.
\iqlookfor{the cumulative-sum trick, an oracle check, and its numerical caveat.}
\end{solution}

\begin{solution}{iq:iv:python:12}
\texttt{[0, 0, 0]}, \texttt{True}, \texttt{[0, 1, 0]}. \texttt{df[df['qty'] < 0]} (a boolean mask) returns a copy, and \texttt{['flag'] = 1}
assigns into that copy, which is discarded: the frame is unchanged and pandas~2.3 emits a \texttt{SettingWithCopyWarning}. The
\texttt{.loc[mask, 'flag'] = 1} form assigns in one step on the frame itself. Under copy-on-write every indexing result behaves as a copy,
so chained assignment never modifies the original and pandas warns about it; \texttt{.loc} remains the way to write.
\iqlookfor{copies from masks, one-step assignment with \texttt{.loc}, and awareness of copy-on-write.}
\end{solution}

\begin{solution}{iq:iv:python:13}
\omcode{code/interviews/24-python/python/iv_py.py}{49}{54}{A frozen dataclass: equality and hash from the
fields.}\label{lst:iv:python:key}
A mutable key changed after insertion keeps its old hash bucket, so lookups by the new value and by the old one can both fail, and
the entry is lost in the dictionary. A frozen dataclass cannot be changed, and its generated \texttt{\_\_hash\_\_} is consistent with
\texttt{\_\_eq\_\_}. A plain tuple works too.
\iqlookfor{hash and equality consistency, immutability, and what breaks otherwise.}
\end{solution}
